% ============================================================================
% Fichier  : partie2/P02_C06_identifier.tex
% Titre    : Identifier les Risques
% Auteur   : MINKA MI NGUIDJOÏ Thierry Emmanuel
% Version  : 1.0 -- novembre 2026
% ============================================================================
% Sources : ISO 31000:2018 (6.4.2) ; IEC 31010:2019 ; ISO/IEC 27005:2022.
% Atelier d'identification de LiZbiZ : cas fictif.
% ============================================================================

\chapter{Identifier les Risques}
\label{chap:identifier}
\chaptermark{Identifier les risques}

\epigraph{%
    « Le risque le plus dangereux\\
    est celui qui ne figure pas dans le registre. »%
}{Principe du cours}

\niveauChapitre{Intermédiaire}{Chapitres~\ref{chap:risque}
à~\ref{chap:cadres}.}

\begin{boxobjectifs}
\begin{itemize}
    \item Formuler un risque sous la forme cause, événement, conséquence.
    \item Choisir et combiner des techniques d'identification.
    \item Classer les risques selon une taxonomie utile.
    \item Partir des objectifs et des processus pour ne rien oublier.
    \item Repérer les signaux faibles et les risques émergents.
\end{itemize}
\end{boxobjectifs}

\bigskip

% ============================================================================
\section*{Question directrice : pourquoi le prestataire ne figurait-il dans aucune liste~?}
% ============================================================================

Lors de l'atelier de mai, chaque directeur de LiZbiZ cite les risques de
son activité. Personne ne cite SAHEL INFOGÉRANCE~: pour la DAF, c'est un
fournisseur~; pour l'informatique, un sous-traitant~; pour les opérations,
un contrat. Un risque qui traverse les frontières des directions échappe
aux listes établies direction par direction. Ce chapitre montre comment
identifier sans laisser de trou.

\medskip
\textbf{Plan du chapitre.}
\begin{enumerate}
    \item Formuler un risque (section~\ref{sec:c06-formuler}).
    \item Partir des objectifs et des processus (section~\ref{sec:c06-processus}).
    \item Les techniques d'identification (section~\ref{sec:c06-techniques}).
    \item Classer~: les taxonomies (section~\ref{sec:c06-taxonomie}).
    \item Les signaux faibles (section~\ref{sec:c06-signaux}).
\end{enumerate}

\separateur

% ============================================================================
\section{Formuler un Risque}
\label{sec:c06-formuler}
% ============================================================================

\begin{boxdef}
\termecle{Énoncé de risque}~: phrase qui relie une \emph{cause} (source ou
vulnérabilité), un \emph{événement} incertain et ses \emph{conséquences}
sur un objectif. Forme recommandée~: « Du fait de [cause], [événement]
peut survenir, ce qui entraînerait [conséquence sur l'objectif]. »
\end{boxdef}

\begin{tableaucours}{L{4.0cm}L{8.2cm}}{%
    \tableauHeader{Énoncé} &
    \tableauHeader{Diagnostic}
}
« Risque cyber » & Une catégorie, pas un risque~: ni cause, ni événement,
ni conséquence \\
« Manque de sauvegardes » & Une cause, pas un risque \\
« Perte de clients » & Une conséquence, pas un risque \\
« Du fait d'un accès distant du prestataire sans second facteur, un
attaquant peut chiffrer les serveurs, ce qui arrêterait la facturation
plusieurs jours » & Un risque bien formulé \\
\end{tableaucours}
\vspace{-0.8em}
\begin{center}{\small\itshape Tableau~6.1. Formuler un risque}\end{center}

\begin{boxalert}
\textbf{Violation typique.} Inscrire au registre des mots-clés (« cyber »,
« réputation », « RH »).

\smallskip
\textbf{Règle~:} un mot-clé ne se cote pas et ne se traite pas. On le
décompose en scénarios formulés, chacun avec sa cause, son événement et sa
conséquence.
\end{boxalert}

% ============================================================================
\section{Partir des Objectifs et des Processus}
\label{sec:c06-processus}
% ============================================================================

L'ISO~31000 rattache le risque aux objectifs~: on identifie d'abord ce que
l'organisation doit atteindre, puis ce qui peut l'en empêcher. Une méthode
sûre consiste à parcourir les \emph{processus} de l'entreprise et, pour
chacun, ses ressources (personnes, systèmes, fournisseurs, locaux) et ses
dépendances.

\begin{tableaucours}{L{3.2cm}L{4.2cm}L{4.8cm}}{%
    \tableauHeader{Processus} &
    \tableauHeader{Ressources critiques} &
    \tableauHeader{Exemple de risque}
}
Facturer & Application intégrée, base de Douala, comptabilité & Arrêt de
l'application plus de deux jours \\
Expédier & Agences, véhicules, suivi des expéditions & Inondation d'un
entrepôt (Bonabéri) \\
Payer & Banque, double validation, trésorerie & Fraude au virement par
usurpation d'un dirigeant \\
Administrer le SI & SAHEL INFOGÉRANCE, comptes à privilèges & Usurpation
du compte du prestataire \\
\end{tableaucours}
\vspace{-0.8em}
\begin{center}{\small\itshape Tableau~6.2. Identifier par les processus}\end{center}

% ============================================================================
\section{Les Techniques d'Identification}
\label{sec:c06-techniques}
% ============================================================================

\begin{tableaucours}{L{3.2cm}L{4.6cm}L{4.4cm}}{%
    \tableauHeader{Technique} &
    \tableauHeader{Apport} &
    \tableauHeader{Limite}
}
Atelier pluridisciplinaire & Croise les points de vue, fait émerger les
risques transversaux & Biais de groupe (Chapitre~\ref{chap:perception}) \\
Entretiens individuels & Libère la parole, détaille un processus & Long~;
vision partielle \\
Listes de contrôle & Ne rien oublier de connu & N'invente rien de nouveau \\
Historique des incidents & Données réelles & Ne montre que le passé \\
Revue de l'écosystème & Prestataires, clients, partenaires & Dépend de
l'information obtenue des tiers \\
Analyse de scénarios & Explore l'inhabituel & Demande une animation
experte \\
\end{tableaucours}
\vspace{-0.8em}
\begin{center}{\small\itshape Tableau~6.3. Techniques d'identification (d'après l'IEC 31010)}\end{center}

\begin{boxprocedure}
\textbf{Conduire un atelier d'identification}
\begin{enumerate}
    \item Envoyer à l'avance les objectifs, la liste des processus et
          l'historique des incidents.
    \item Faire écrire à chacun, en silence, ses risques pendant dix
          minutes, avant toute discussion.
    \item Regrouper, reformuler selon la forme cause, événement,
          conséquence~; supprimer les doublons.
    \item Parcourir chaque processus et chaque tiers~: « qu'a-t-on
          oublié~? »
    \item Désigner pour chaque risque un propriétaire provisoire.
    \item Consigner les risques écartés et la raison de leur exclusion.
\end{enumerate}
\end{boxprocedure}

\begin{boxscenario}{Opération LIZBIZ. L'atelier de mai}
\textbf{Constatation.} L'atelier réunit huit personnes de six directions.
Après la phase d'écriture silencieuse, 63~propositions~; après
regroupement et reformulation, 41~risques. Le parcours des tiers en ajoute
cinq, dont l'usurpation du compte du prestataire et la défaillance de
l'unique opérateur de télécommunications de Garoua.

\medskip
\textbf{Analyse} (confiance élevée). Les cinq risques ajoutés par la revue
de l'écosystème n'avaient été cités par personne~: l'identification par
direction laisse passer les risques qui n'appartiennent à aucune.
\end{boxscenario}

% ============================================================================
\section{Classer : les Taxonomies}
\label{sec:c06-taxonomie}
% ============================================================================

Une \termecle{taxonomie des risques} répartit les risques en familles
stables, pour vérifier que chacune a été examinée et consolider les
résultats. Elle ne doit être ni trop fine (on ne s'y retrouve plus), ni
trop grossière (« opérationnel » couvre tout).

\begin{tableaucours}{L{3.0cm}L{9.2cm}}{%
    \tableauHeader{Famille} &
    \tableauHeader{Exemples chez LiZbiZ}
}
Stratégique & Perte d'un client majeur~; arrivée d'un concurrent \\
Opérationnel & Inondation d'entrepôt~; panne de l'application~; erreur
de saisie \\
Numérique & Rançongiciel~; fuite de données~; usurpation d'accès \\
Financier & Fraude au virement~; impayé d'un client \\
Conformité & Méconnaissance de l'article~26~; données personnelles \\
Tiers & Défaillance du prestataire~; dépendance à un opérateur \\
\end{tableaucours}
\vspace{-0.8em}
\begin{center}{\small\itshape Tableau~6.4. Une taxonomie pour LiZbiZ}\end{center}

% ============================================================================
\section{Les Signaux Faibles}
\label{sec:c06-signaux}
% ============================================================================

\begin{boxdef}
\termecle{Signal faible}~: information partielle, isolée ou ambiguë, qui
peut annoncer un risque nouveau ou l'aggravation d'un risque connu.
\end{boxdef}

Les signaux faibles se trouvent dans les quasi-incidents, les plaintes des
clients, les alertes des autorités, l'actualité des pairs. En janvier
2026, une entreprise de logistique voisine avait été frappée par un
rançongiciel entré par l'accès d'un prestataire~; l'information avait
circulé, sans être rapprochée de la situation de LiZbiZ.

\begin{boxobjet}{R-03 · Défaillance de l'opérateur de Garoua}
\textbf{Scénario~:} du fait d'un opérateur unique à Garoua, une coupure
prolongée du réseau isole l'agence, ce qui empêche le suivi et la
facturation des expéditions du Nord. \textbf{Cotation~:} vraisemblance 2,
impact 3. \textbf{Traitement~:} réduire (second lien par un autre
opérateur, mode dégradé papier). \textbf{Propriétaire~:} M.~NGUELE.
\textbf{Indicateur~:} durée cumulée des coupures du trimestre.
\end{boxobjet}

\begin{boxmetier}{Gestionnaire de risques}
\textbf{Tenir une veille.} Mme~MOUAFO tient une liste des incidents
publiés dans le secteur et, pour chacun, se demande~: « cela pourrait-il
nous arriver, par où, et l'avons-nous au registre~? »
\end{boxmetier}

\separateur

\begin{boxresume}
\textbf{Ce qu'il faut retenir de ce chapitre~:}
\begin{itemize}
    \item \textbf{Un risque se formule}~: cause, événement, conséquence sur
          un objectif.
    \item \textbf{On part des objectifs et des processus}, et l'on parcourt
          les tiers.
    \item \textbf{Les techniques se combinent}~: atelier, entretiens,
          listes, historique, écosystème, scénarios.
    \item \textbf{Une taxonomie} vérifie la couverture et consolide.
    \item \textbf{Les signaux faibles} nourrissent le registre~; la veille
          les rapproche de sa propre situation.
\end{itemize}
\end{boxresume}

\begin{boxinfo}
\textbf{Points de contrôle de l'élève.} Je sais~: formuler un risque~;
identifier par les processus~; animer un atelier~; choisir une technique~;
classer selon une taxonomie~; exploiter un signal faible.
\end{boxinfo}

\bigskip

\begin{boxexo}[title={\ding{45}\quad Exercice~6.1. Diagnostiquer \niveaubase}]
Dites s'il s'agit d'une cause, d'un événement, d'une conséquence ou d'un
risque formulé~: « absence de second facteur »~; « publication des données
clients »~; « arrêt de la facturation ».
\end{boxexo}

\begin{boxexo}[title={\ding{45}\quad Exercice~6.2. Formuler \niveaubase}]
Formulez trois risques pour le processus « Payer » de LiZbiZ.
\end{boxexo}

\begin{boxexo}[title={\ding{45}\quad Exercice~6.3. Classer \niveaubase}]
Classez dans la taxonomie du Tableau~6.4 les risques R-01, R-02 et R-03.
\end{boxexo}

\begin{boxexo}[title={\ding{45}\quad Exercice~6.4. L'atelier \niveaumoyen}]
Pourquoi faire écrire chacun en silence avant la discussion~? Quel biais
du Chapitre~\ref{chap:perception} la règle combat-elle~?
\end{boxexo}

\begin{boxexo}[title={\ding{45}\quad Exercice~6.5. Les tiers \niveaumoyen}]
Dressez la liste des tiers dont dépend la facturation de LiZbiZ, et un
risque pour chacun.
\end{boxexo}

\begin{boxexo}[title={\ding{45}\quad Exercice~6.6. Écarter \niveaumoyen}]
Pourquoi consigner les risques écartés et la raison de leur exclusion~?
\end{boxexo}

\begin{boxexo}[title={\ding{45}\quad Exercice~6.7. Signal faible \niveaumoyen}]
L'attaque d'une entreprise voisine en janvier était un signal faible.
Décrivez le circuit qui aurait permis de le rapprocher de LiZbiZ.
\end{boxexo}

\begin{boxexo}[title={\ding{45}\quad Exercice~6.8. Inondation \niveaumoyen}]
Identifiez les risques de l'entrepôt de Bonabéri pendant la saison des
pluies~: causes, événements, conséquences.
\end{boxexo}

\begin{boxexo}[title={\ding{45}\quad Exercice~6.9. Une banque \niveauexpert}]
Préparez l'atelier d'identification des risques du système central de la
BANQUE SANAGA~: participants, documents, déroulé, taxonomie.
\end{boxexo}

\begin{boxexo}[title={\ding{45}\quad Exercice~6.10. Complétude \niveauexpert}]
Comment vérifier qu'un registre de 46~risques est complet~? Proposez trois
contrôles indépendants.
\end{boxexo}

\bigskip

\begin{boxinfo}
\textbf{Pour aller plus loin.}
\begin{itemize}
    \item ISO~31000~\cite{iso31000_2018}~; IEC~31010~\cite{iec31010}.
    \item ISO/IEC~27005~\cite{iso27005_2022}.
\end{itemize}
\end{boxinfo}

\begin{boxinfo}
\textbf{Transition.} Une fois identifiés, les risques doivent être cotés.
Le Chapitre~\ref{chap:matrices} présente les échelles et les matrices,
avec leurs limites.
\end{boxinfo}
